\begin{abstract}
The distributed surface mass balance (SMB) of mountain glaciers is a key forcing for glacier evolution and subsequent meltwater availability projections. Yet, direct in situ observations cover 
%GJ: a tiny portion of (it is probably much less than 0.2, depends on how you define it, spatial/temporal coverage ...), "a tiny portion" is clear enough for an abstract
less than 0.2\,\% 

of the world's glacierised area. Here, we present a new SMB inference method, implemented within the open-source Instructed Glacier Model (IGM), that recasts the problem as 
%GJ: "as a inverse problem." you may keep it general statement here since you develop the differentiability right after.
a fully differentiable transient inversion. 
The inversion scheme uses observed surface velocities to initialize a dynamically consistent glacier model, then inverts multi-temporal surface elevation changes (DEMs) to estimate the spatial distribution of surface mass balance that best explains the observed glacier evolution.  We apply the method to three glaciers in the Alps---Argentière, Great Aletsch and Rhône---and validate the inferred profiles against in situ stake measurements. A sensitivity test shows that the surface-velocity misfit alone bounds the inferred SMB on stake-free glaciers. Our method is generic, computationally efficient and requires no glacier-specific tuning, making it readily transferable to any glacier with accurate multi-temporal DEMs and surface velocities, opening a path toward regional-scale, continuous SMB inferences.
\end{abstract}